\section{Introduction \& Motivation}

\begin{frame}[allowframebreaks]{Why Audio in NLP?}
  \large{Speech is the \textbf{most natural form of communication}.}

  \vspace{0.5em}

  \textbf{What:} Where speech interfaces are essential.
  \begin{itemize}
    \item Voice assistants (e.g., Amazon Echo)
    \item Accessibility: captions, screen readers
    \item Hands-free operation (cars, IoT)
    \item Multilingual interaction & translation
  \end{itemize}
  \vspace{0.5em}
  \textbf{Why:} Speech enables more intuitive and inclusive interaction.

\framebreak
  
  \fetchconvertimage{https://i.gzn.jp/img/2017/03/05/alexa-siri-cortana-all-female/00.png}{images/audio-nlp-eval/echo_dot.png}{width=1.05\linewidth,height=0.9\textheight,keepaspectratio}
\end{frame}

\begin{frame}{Motivation}
    \begin{itemize}
        \setlength{\itemsep}{1em}
        \item Audio data is abundant and diverse, encompassing speech, music, and environmental sounds.
        \item Traditional NLP methods focus on text, but audio data requires specialized techniques for processing.
        \item Many NLP systems now start from speech, not just text (e.g., live captions, real-time translation).
        \item The rise of voice assistants and audio-based applications highlights the need for effective audio processing in NLP.
        \item Audio signals contain rich information that can enhance understanding and interaction in various applications.
    \end{itemize}
\end{frame}
